\section{思考与总结}

\subsection{开发过程中遇到的问题与解决方案}

\subsubsection{多源数据的标签一致性}

课程评价来自公开评论与人工构造场景。
不同来源的表达风格和信息密度存在差异。
项目首先统一积极、中性、消极的标注标准；
再对三类标签进行均衡采样，并固定训练、验证、测试划分。
这样可以让传统模型、BERT 和混合流程在相同数据基础上比较；
也避免数据来源差异直接干扰实验结论。

\subsubsection{中性与混合评价边界模糊}

三分类任务中，中性类别不仅包括没有明显态度的文本；
也包括建议型表达和正负混合评价。
传统 TF-IDF 模型容易只关注强情感词；
BERT 也可能将转折句推向某一极性。
项目通过统一标签定义、转折规则和正负信号平衡判断，
将“讲解清楚但作业太多”归入中性，并保留各课程维度的具体问题。

\subsubsection{非规范英文表达影响模型输入}

在线评论中存在缩写、漏字、重复字母和拼写错误。
若全面启用通用拼写纠错，容易误改课程名和技术词；
若完全不处理，又会损失情感信号。
最终方案是在情感通道中使用受限词典和极小编辑操作，只纠正高置信度情感词。
原文仍用于证据展示，避免证据与学生原话不一致。

\subsubsection{长文本截断导致信息丢失}

BERT 的输入长度有限，直接截断可能只保留评价开头。
项目实现带重叠的滑动窗口，并按有效 Token 数加权聚合各窗口概率。
同时限制最大窗口数并记录截断元数据。
这样可以在计算量可控的前提下，尽量保留长评价中的后续信息。

\subsubsection{模型与外部服务可能不可用}

BERT 权重体积较大，加载需要时间。
课堂演示时，ChatECNU API 也可能因为网络或配置问题不可用。
项目将模型加载设计为懒加载和实例复用；
情感分析提供 BERT、TF-IDF、规则三级后端；
建议生成提供 API 与本地模板两级方案。
前端显示实际运行状态，避免用户误以为系统始终在使用同一种模型。

\subsubsection{固定案例不能替代泛化评估}

固定案例适合演示，但规则可能因为覆盖这些表达而取得较高成绩。
项目额外建立 48 条独立压力测试，并按表达类型和语言统计结果。
压力测试暴露了讽刺和隐含抱怨的不足；
也避免将 12/12 的演示结果解释为整体准确率 100\%。

\subsection{局限性与使用边界}

\begin{enumerate}
  \item \textbf{中文真实数据不足。}当前主要量化评估基于英文 Coursera 评论；
        120 条中文样本只用于流程验证和课堂演示，不能支撑完整中文泛化结论。
  \item \textbf{真实场景覆盖仍可扩展。}当前数据尚未覆盖更多学科、学期和院校环境；
        跨场景泛化能力仍需继续验证。
  \item \textbf{课程维度依赖词典。}领域词典易解释、部署轻量；
        但难以识别没有显式关键词的隐含主题，也需要人工维护新课程术语。
  \item \textbf{讽刺和隐含抱怨识别较弱。}压力测试中讽刺通过率仅 1/6；
        说明有限规则和当前训练数据不足以处理复杂语用现象。
  \item \textbf{大模型建议需要人工复核。}系统要求 LLM 引用证据；
        但生成内容仍不能直接作为教学决策，教师需要结合课程实际情况判断。
  \item \textbf{BERT 模型文件较大。}BERT 权重约 711 MB；
        在 CPU 上可以运行，但首次加载和推理速度不如传统模型。
  \item \textbf{尚未进行真实用户研究。}当前主要验证算法和系统流程；
        还缺少教师实际使用后的可用性评价与改进效果追踪。
\end{enumerate}

\subsection{未来改进方向}

\begin{itemize}
  \item 与真实课程合作，在脱敏和授权前提下收集跨课程、跨学期评价，并由多人复核标签；
  \item 对不同语言和不同课程场景进行分层评估，增加置信区间和跨学期测试；
  \item 将课程维度识别升级为多标签分类，同时继续保留证据抽取和规则解释；
  \item 针对讽刺、隐含抱怨和复杂否定构建专门训练集，而不是无限增加手写规则；
  \item 引入模型蒸馏、量化或 ONNX Runtime，减小部署体积并提升 CPU 推理速度；
  \item 加入课程、学期和教师维度的历史趋势对比，以及人工反馈修正机制；
  \item 对 LLM 建议增加事实一致性检查、敏感信息脱敏和人工确认工作流。
\end{itemize}

\subsection{个人收获与反思}

通过这次项目，我们对自然语言处理系统的完整开发流程有了更具体的认识。
最开始我们更多关注模型能否完成情感分类，但随着项目推进，才发现一个可用的系统还需要考虑数据来源、标签定义、输入清洗、模型回退、界面展示和结果解释等多个环节。
尤其是在课程评价这样的实际场景中，系统不能只给出一个情感标签，还需要告诉用户判断依据是什么、问题集中在哪些课程环节，以及这些结果是否值得进一步人工复核。

在模型实现过程中，我们对传统机器学习方法和 BERT 模型的差异有了更直观的理解。
TF-IDF Logistic Regression 训练和部署都比较轻量，适合作为稳定的基础模型；
multilingual BERT 在处理上下文、转折和否定表达时效果更好，但也带来了模型文件较大、加载较慢和环境依赖更多的问题。
因此，项目最后没有简单地只保留一个模型，而是设计了 BERT、传统模型和规则方法的分层回退流程。
这个过程让我们意识到，课程项目中的“效果好”不仅指指标更高，也包括在不同运行环境下能否稳定工作。

项目中比较有启发的一点是标签和评价场景之间的关系。
课程评论里经常出现“老师讲得很清楚，但是作业太多”这类混合表达，如果只按积极或消极二分，很容易丢失学生真正想表达的问题。
因此，我们在情感分类之外加入了课程维度、关键词和原文证据，希望把评论中的具体关注点保留下来。
这样做虽然不能完全解决语义理解问题，但能让教师更快看到学生反馈集中在哪些方面，也便于判断系统分析是否合理。

在实验评估方面，我们也更加注意不同结果的使用边界。
固定案例适合课堂展示和流程检查，但不能代表模型的泛化能力；独立测试集更适合用于模型对比；压力测试则能暴露系统在复杂表达下的不足。
比如当前系统对缩写、拼写错误、中英混合和长文本有一定改进，但面对讽刺、隐含抱怨等表达时仍然不稳定。
这些结果提醒我们，报告中不能只展示好看的指标，也需要说明失败情况和适用范围。

这次项目也暴露了一些不足。
由于时间和数据限制，中文真实课程评价样本仍然偏少，中文和中英混合输入更多依靠规则和小规模测试验证；
课程维度识别主要基于词典，对隐含主题的识别能力有限；
LLM 生成的教学建议虽然能提高可读性，但仍然需要教师根据真实教学情境判断。
后续如果继续完善，可以重点补充真实中文评价数据，改进维度标注方式，并对复杂表达进行更系统的测试。

总体来说，这次项目让我们不只是练习了文本预处理、TF-IDF、BERT 和模型评估等课程内容，也尝试把这些方法整合成一个面向具体任务的应用系统。
相比单独完成一个分类模型，我们对 NLP 项目的工程实现、结果解释和使用边界都有了更实际的理解。

\subsection{总结}

CourseInsight 已实现从中文、英文及中英混合课程评价输入到结构化教学反馈输出的基本闭环。
系统包含双语预处理、BERT 情感分类、规则校正、模型回退、课程维度证据、关键词、建议生成、
批量分析、固定案例验证和模型评估页面。
实验结果说明，混合流程在固定案例和独立压力测试中均优于单一模型方案；
同时也暴露了讽刺识别、隐含抱怨和中文真实数据不足等问题。

作为课程项目，项目具备完整源代码、部署说明、模型指标、消融实验、独立压力测试和可运行界面。
作为应用原型，它可以把大量课程评价快速整理为可复核的改进线索。
后续通过补充真实场景数据、升级课程维度模型和开展教师用户测试，可以进一步提高可靠性。
